\title{DeepSeekMath: Pushing the Limits of Mathematical Reasoning in Open Language Models}

\begin{abstract}
The inherently complex and highly structured character of mathematical reasoning presents notable difficulties for language models. In this work, we present DeepSeekMath~7B, an extension of the DeepSeek-Coder-Base-v1.5 7B model, further pre-trained on 120B tokens related to mathematics, curated from Common Crawl along with supplementary natural language and programming code datasets. Without the use of external tools or ensemble voting strategies, DeepSeekMath~7B attains an impressive 51.7\% score on the competition-caliber MATH benchmark, coming close to the performance displayed by models such as Gemini-Ultra and GPT-4. Employing self-consistency with 64 sampled solutions, the model achieves a score of 60.9\% on MATH. The advanced mathematical reasoning skills demonstrated by DeepSeekMath~are primarily a result of two contributions: Firstly, the development and deployment of a rigorous data filtering and selection system that leverages the breadth of public web resources; secondly, the design of Group Relative Policy Optimization (GRPO), an adaptation of Proximal Policy Optimization (PPO), which both strengthens the model’s mathematical reasoning and improves the memory efficiency associated with PPO.
\end{abstract}

\section{Introduction}

Large language models (LLMs) have fundamentally transformed methodologies for mathematical reasoning within the field of artificial intelligence, leading to marked progress on quantitative reasoning benchmarks \citep{MATH} as well as on geometry reasoning tasks \citep{trinh2024solving}. Additionally, these models have demonstrated substantial value in supporting human users tackling challenging mathematical questions \citep{tao}. Despite these achievements, leading models like GPT-4 \citep{gpt4} and Gemini-Ultra \citep{gemini} remain inaccessible to the public, while available open-source alternatives exhibit significantly inferior performance.

This paper presents DeepSeekMath, a specialized language model designed for mathematics that substantially surpasses the abilities of existing open-source systems and nearly matches GPT-4 on standard academic evaluations. Central to this advancement is the creation of the DeepSeekMath Corpus—a comprehensive, high-caliber pre-training dataset containing 120B math-related tokens. The dataset is assembled from the Common Crawl (CC) using a classifier built on the fastText framework \citep{joulin2016fasttext}. Initially, this classifier is trained with positive examples drawn from OpenWebMath \citep{openwebmath} and a representative set of unrelated web content as negative examples. Using this version of the classifier, further positive examples are mined from the CC, which are then subjected to human annotation for quality control. The classifier is iteratively improved by retraining on this augmented and refined dataset. Evaluation demonstrates that the resulting corpus is of notable quality: our foundational model, DeepSeekMath-Base 7B, achieves scores of 64.2\% on GSM8K \citep{gsm8k} and 36.2\% on the challenging MATH competition dataset \citep{MATH}, outperforming models such as Minerva 540B \citep{minerva}. Furthermore, the DeepSeekMath Corpus is constructed to be multilingual, yielding enhanced performance on Chinese mathematical benchmarks \citep{wei2023cmath,agieval}. We propose that our methodology for compiling and processing mathematical data offers a valuable baseline for future research and envision significant opportunities for further advancement.

DeepSeekMath-Base is initially constructed using DeepSeek-Coder-Base-v1.5 7B \citep{deepseek-coder}, as our findings suggest that beginning with a code-centric model outperforms initializing from a generic LLM. Additionally, we have found that the specialized math training enhances the model's performance on benchmarks such as MMLU \citep{mmlu} and BBH \citep{bbh}, revealing that these improvements extend beyond mathematics to bolster the model’s general reasoning skills.

Following the pre-training phase, DeepSeekMath-Base undergoes mathematical instruction tuning that incorporates chain-of-thought \citep{cot}, program-of-thought \citep{pot,pal}, and tool-integrated reasoning \citep{tora} data. This produces DeepSeekMath-Instruct 7B, which outperforms all existing 7B models in its class and demonstrates capabilities on par with much larger 70B open-source instruction-tuned models.

In addition, we propose Group Relative Policy Optimization (GRPO), a new reinforcement learning (RL) approach derived from Proximal Policy Optimization (PPO) \citep{schulman2017proximal}. Unlike PPO, GRPO eliminates the need for a critic model by using group-based scores to compute the baseline, thereby substantially lowering computational overhead during training. Leveraging only a subset of English instruction tuning datasets, GRPO achieves notable gains over the already strong DeepSeekMath-Instruct baseline, with marked improvements on both in-domain benchmarks (GSM8K: 82.9\% $\rightarrow$ 88.2\%, MATH: 46.8\% $\rightarrow$ 51.7\%) and on out-of-domain challenges (for instance, CMATH: 84.6\% $\rightarrow$ 88.8\%) throughout the RL process. Furthermore, we present a unified framework to systematically analyze and interpret various approaches—including Rejection Sampling Fine-Tuning (RFT) \citep{yuan2023scaling}, Direct Preference Optimization (DPO) \citep{dpo}, PPO, and GRPO—demonstrating that they all can be viewed as instances of direct or simplified RL paradigms. To thoroughly examine the key components of this framework, we conduct a comprehensive suite of experiments, such as comparisons between online and offline training, outcome-based versus process-based supervision, and single-turn versus iterative RL. Finally, we elucidate the mechanisms through which our RL approach enhances instruction-tuned model performance and outline future research avenues for developing even more robust RL techniques within this unified framework.

\subsection{Contributions}
Our contribution includes scalable math pre-training, along with the exploration and analysis of reinforcement learning.

\textbf{Math Pre-Training at Scale}

\begin{itemize}[topsep=0pt]
    \item This study demonstrates that the openly available Common Crawl dataset is a rich resource for mathematical applications. Utilizing a carefully crafted data selection process, we assembled the DeepSeekMath Corpus—a high-caliber dataset comprising 120B tokens derived from web pages specifically filtered for mathematical relevance. This corpus is nearly 7 times larger than the collection of math web pages in Minerva \citep{minerva}, and exceeds the recently launched OpenWebMath \citep{openwebmath} by a factor of 9.

\item DeepSeekMath-Base 7B, our foundational model, matches the performance level of Minerva 540B \citep{minerva}. This suggests that sheer model size is not the sole determinant of effectiveness in mathematical reasoning. Notably, robust outcomes can also be attained by smaller models when they are pre-trained on high-caliber data.

\item We present insights gained from our experiments involving math training. Preparing models through code training before engaging them in math tasks leads to enhanced mathematical problem-solving, regardless of whether external tools are employed. This result provides some evidence addressing the ongoing question: \textit{does code training enhance reasoning skills?} Our results indicate that, at least in the context of mathematical reasoning, it does.

    \item Although training on arXiv papers is common, especially in many math-related papers, it brings no notable improvements on all mathematical benchmarks adopted in this paper.
\end{itemize}

\textbf{Exploration and Analysis of Reinforcement Learning}

\begin{itemize}[topsep=0pt]
    \item We present Group Relative Policy Optimization (GRPO), a resource-efficient and robust reinforcement learning method. Unlike conventional approaches such as Proximal Policy Optimization (PPO), GRPO eliminates the need for a critic network by utilizing group-based score estimates for the baseline, resulting in considerably lower computational demands.
    \item Through empirical evaluation, we show that GRPO substantially improves the capability of our instruction-tuned model, DeepSeekMath-Instruct, relying exclusively on instruction-tuning datasets. Notably, we also detect better generalization on out-of-domain tasks throughout reinforcement learning.
    \item We establish a comprehensive framework that encompasses a range of existing techniques, including RFT, DPO, PPO, and GRPO. Through a series of thorough experiments—such as comparisons between online and offline learning, outcome-based and process-based supervision, as well as single-turn versus iterative reinforcement learning—we systematically analyze the foundational aspects of this framework.
    \item Drawing on insights from our unified perspective, we investigate the key factors underlying reinforcement learning’s effectiveness and identify several promising future research directions for enhancing reinforcement learning for large language models (LLMs).
\end{itemize}

\subsection{Summary of Evaluations and Metrics}

\begin{itemize}[topsep=0pt]
    \item \textbf{English and Chinese Mathematical Reasoning}: Our evaluation covers both English and Chinese datasets, spanning a wide range of mathematical tasks from elementary to collegiate difficulty. For English, we utilize benchmarks such as GSM8K \citep{gsm8k}, MATH \citep{MATH}, SAT \citep{llemma}, OCW Courses \citep{minerva}, and MMLU-STEM \citep{mmlu}. The Chinese evaluation includes MGSM-zh \citep{mgsm}, CMATH \citep{wei2023cmath}, Gaokao-MathCloze \citep{agieval}, and Gaokao-MathQA \citep{agieval}. The assessment examines the capacity of the models to produce comprehensive textual solutions independently of external tools, as well as their proficiency in problem-solving through Python.

On English-language evaluation tasks, DeepSeekMath-Base demonstrates performance on par with the proprietary Minerva 540B model \citep{minerva}, and consistently outperforms all available open-source base models, including Mistral 7B \citep{mistral} and Llemma-34B \citep{llemma}, irrespective of whether these alternatives incorporate mathematical pre-training, frequently by a wide margin. Importantly, DeepSeekMath-Base achieves even stronger results on Chinese evaluation benchmarks. This can likely be attributed to our deviation from previous approaches \citep{minerva,llemma} that relied exclusively on English-language mathematical pre-training data; instead, we augment the dataset with high-quality materials from additional languages. By further applying mathematical instruction tuning and reinforcement learning, the resultant DeepSeekMath-Instruct and DeepSeekMath-RL models achieve outstanding outcomes, with both models surpassing 50\% accuracy on the competition-grade MATH dataset—marking a first among open-source solutions.

\item \textbf{Formal Mathematics}: DeepSeekMath-Base is assessed on the task of transforming informal statements into formalized proofs, utilizing the dataset from \citep{dsp_proof} based on miniF2F \citep{minif2f} with Isabelle \citep{isabelle} serving as the proof assistant. The model exhibits notable few-shot performance in automatic formalization tasks.
\item \textbf{Natural Language Understanding, Reasoning, and Code}: To thoroughly measure the models’ capacities in general comprehension, logical reasoning, and programming, we evaluate DeepSeekMath-Base on several benchmarks: Massive Multitask Language Understanding (MMLU) \citep{mmlu}, which features 57 diverse multiple-choice subject areas; BIG-Bench Hard (BBH) \citep{bbh}, comprising 23 intricate problems necessitating multi-step reasoning; as well as HumanEval \citep{codex} and MBPP \citep{mbpp}, both established for code model assessment. Pre-training on mathematical tasks enhances performance not just in mathematics, but also in linguistic understanding and reasoning challenges.

\section{Math Pre-Training}

\subsection{Data Collection and Decontamination} This section describes the methodology used to build the DeepSeekMath Corpus utilizing data sourced from Common Crawl. As illustrated in Figure \ref{fig:data_collect}, we introduce a step-by-step pipeline designed to efficiently assemble an extensive mathematical corpus from Common Crawl, initiating the process with a high-quality seed dataset comprising a limited set of mathematics-focused data. Importantly, the outlined procedure is generalizable and may also be leveraged for the development of datasets in other specialized fields, like programming.

To begin, OpenWebMath \citep{openwebmath}, which comprises a curated set of high-quality mathematical web content, serves as our foundational seed corpus. Leveraging this collection, we build a fastText model \citep{joulin2016fasttext} aimed at retrieving additional web pages that are similar in nature to OpenWebMath. We draw 500,000 entries at random from the seed corpus to act as positive examples and supplement these with another 500,000 randomly selected web pages from Common Crawl as negative examples. Training is carried out using a publicly available implementation\footnote{\url{https://fasttext.cc}}, with the following parameters: 256-dimensional vectors, a learning rate of 0.1, a maximum word n-gram length of 3, a minimum occurrence threshold of 3 per word, and 3 epochs of training. In order to manage the size of the initial Common Crawl dataset, we apply URL-based deduplication as well as near-duplicate filtering, ultimately yielding a set of 40 billion HTML web pages. The trained fastText model is then used to identify mathematical web pages from this deduplicated dataset. To ensure the quality of the resulting mathematical content, we rank the retrieved pages using the scores output by the fastText model and retain only those with the highest ranks. The final selection volume is further refined by conducting pre-training trials on datasets of increasing size, namely the top 40B, 80B, 120B, and 160B tokens. For the first phase, we opt to retain the subset comprising the top 40B tokens.

Following the initial round of data acquisition, a significant portion of mathematical web pages remains undiscovered. This shortfall is primarily attributed to the limited diversity within the set of positive examples used to train the fastText model. To enhance the model, we expand our selection of mathematical web sources, thereby broadening the seed corpus. Our strategy involves segmenting the entire Common Crawl by distinct domains, with a domain defined as a collection of web pages sharing a common base URL. For each domain, we compute the proportion of pages successfully collected during the first iteration. If more than 10\% of a domain’s pages have been retrieved, the domain is considered math-oriented (e.g., \url{mathoverflow.net}).

We then perform a manual review of the URLs within these selected domains that contain mathematical materials (such as \url{mathoverflow.net/questions}). Previously uncollected pages associated with these URLs are incorporated into the seed corpus. By doing so, we obtain a richer set of positive examples for training, which allows for the development of a more robust fastText model, capable of recovering a higher volume of mathematical content in the next cycle. After four data collection rounds, our method yields 35.5M mathematical web pages, summing up to a total of 120B tokens. In the fourth round, it becomes apparent that almost 98\% of the content had already been captured in the preceding iteration, leading us to discontinue further data collection.

In order to prevent contamination of benchmarks, we adopt the filtering approach described in \cite{deepseek-coder}: all web pages that include questions or answers from established English mathematical benchmarks, such as GSM8K \citep{gsm8k} and MATH \citep{MATH}, as well as Chinese resources like CMATH \citep{wei2023cmath} and AGIEval \citep{agieval}, are excluded. Specifically, our criterion stipulates that any passage containing a sequence of 10 consecutive tokens exactly identical to any segment from these evaluation benchmarks will be eliminated from the math training dataset. For those benchmark texts with fewer than 10 but at least 3 tokens, precise matching is likewise employed to excise any compromised web content.

\subsection{Validating the Quality of the DeepSeekMath~Corpus}

We run pre-training experiments to investigate how the DeepSeekMath~Corpus is compared with the recently released math-training corpora:

\begin{itemize}[topsep=0pt]
    \item \textbf{MathPile} \citep{mathpile}: an extensive dataset containing 8.9B tokens, assembled from various sources such as textbooks, Wikipedia, ProofWiki, CommonCrawl, StackExchange, and arXiv. Notably, arXiv accounts for more than 85\% of this compilation;
    \item \textbf{OpenWebMath} \citep{openwebmath}: a 13.6B-token corpus comprised of mathematical text identified and extracted from the CommonCrawl dataset;
    \item \textbf{Proof-Pile-2} \citep{llemma}: this dataset integrates material from OpenWebMath, 10.3B tokens of code from AlgebraicStack, and 28.0B tokens from arXiv academic papers. For Proof-Pile-2 evaluations, we adopt the arXiv:Web:Code distribution of 2:4:1 as described in \cite{llemma}.
\end{itemize}

\subsubsection{Training Setting}
\label{sec:quality-policy}
We perform specialized math-focused training on a foundational language model comprising 1.3B parameters, architecturally aligned with DeepSeek LLMs \citep{deepseek-llm}, referred to as DeepSeek-LLM 1.3B. Each mathematical dataset is utilized to train a distinct model for 150B tokens. All model training leverages the lightweight and resource-efficient HAI-LLM framework \citep{haillm}. Adopting the training regimen established for DeepSeek LLMs, we employ the AdamW optimizer \citep{adamW} configured with $\beta_1=0.9$, $\beta_2=0.95$, and $\mathrm{weight\_decay}=0.1$. The learning rate follows a multi-phase schedule: it peaks following 2,000 steps of linear warmup, then decays to 31.6\% of the peak value at the 80\% point of the training, and finally reduces further to 10.0\% of the maximum after 90\% completion. The peak learning rate is fixed at 5.3e-4. Training operates with a batch size of 4M tokens and a context window of 4K tokens.

\subsubsection{Evaluation Results}

\textbf{The DeepSeekMath~Corpus is of high quality, covers multilingual mathematical content, and is the largest in size.}

\begin{itemize}[topsep=0pt]
    \item \textbf{High-quality}:
    We assess the effectiveness of downstream tasks across 8 mathematical benchmarks using few-shot chain-of-thought prompting \cite{cot}. Table \ref{tab:corpora_comparison} clearly indicates that models trained with the DeepSeekMath~Corpus consistently outperform others. Furthermore, as illustrated in Figure \ref{fig:corpora_comparisons}, a model trained with the DeepSeekMath Corpus exhibits superior results over Proof-Pile-2 after processing 50B tokens (which equates to one complete epoch of Proof-Pile-2), suggesting that the typical quality of DeepSeekMath Corpus data is notably higher.
    \item \textbf{Multilingual}:
    The DeepSeekMath~Corpus provides multilingual data, with English and Chinese making up the majority of its contents. Table \ref{tab:corpora_comparison} demonstrates that leveraging the DeepSeekMath~Corpus in training boosts mathematical reasoning for both English and Chinese, unlike most available mathematical corpora, which are dominated by English-language data and offer only marginal—or even negative—gains in Chinese mathematical reasoning.
    \item \textbf{Large-scale}:
    In terms of scale, the DeepSeekMath~Corpus surpasses prior mathematical corpora by a substantial margin. Figure \ref{fig:corpora_comparisons} shows that the DeepSeek-LLM 1.3B model, when exposed to the DeepSeekMath~Corpus, not only learns faster but also achieves more sustained improvements over time. Conversely, baseline datasets are considerably smaller and are cycled through multiple times during training, which leads to models hitting a performance ceiling rapidly.
\end{itemize}

\subsection{Training and Evaluating DeepSeekMath-Base 7B}

This section presents DeepSeekMath-Base 7B, a foundational model demonstrating advanced mathematical reasoning skills. We begin with DeepSeek-Coder-Base-v1.5 7B \citep{deepseek-coder} as the initial checkpoint and train it over 500B tokens. The dataset used in training is distributed as follows: 56\% comes from the DeepSeekMath Corpus, 4\% originates from AlgebraicStack, 10\% is sourced from arXiv, 20\% consists of Github code, while the final 10\% includes natural language samples from Common Crawl in English and Chinese. We employ the training procedures outlined in Section \ref{sec:quality-policy}, with two exceptions: the peak learning rate is increased to 4.2e-4, and the batch size is adjusted to 10M tokens.

Our study undertakes an extensive evaluation of DeepSeekMath-Base 7B's mathematical proficiency, emphasizing its performance in generating independent mathematical solutions unaided by external resources, utilizing tools for mathematical problem solving, and executing formal theorem proofs. In addition to mathematics, we furnish a broader analysis of the foundational model, covering its aptitude in natural language comprehension, logical reasoning, and coding abilities.

\paragraph{Mathematical Problem Solving with Step-by-Step Reasoning}

We assess the problem-solving capabilities of DeepSeekMath-Base by employing few-shot chain-of-thought prompting \citep{cot} over eight distinct benchmarks in both English and Chinese. These benchmarks include datasets focused on quantitative reasoning, such as GSM8K \citep{gsm8k}, MATH \citep{MATH}, and CMATH \citep{wei2023cmath}, as well as multiple-choice formats, including MMLU-STEM \citep{mmlu} and Gaokao-MathQA \citep{agieval}. The selected tasks span a broad array of mathematical disciplines, ranging from elementary to university-level challenges.

According to Table \ref{tab:base_math_cot}, DeepSeekMath-Base 7B consistently achieves the highest scores across all eight evaluated benchmarks when compared to other open-source base models, such as the commonly adopted general model Mistral 7B \citep{mistral} and the more recent Llemma 34B \citep{llemma}, the latter of which was trained on Proof-Pile-2 \citep{llemma} for mathematical tasks. In particular, on the challenging MATH benchmark designed for competition-level problems, DeepSeekMath-Base outperforms other available open-source models by more than 10 percentage points. Furthermore, it exceeds the performance of Minerva 540B \citep{minerva}, a much larger closed-source model based on PaLM \citep{palm} with additional math-specific training, despite Minerva's substantially greater scale—at 77 times the parameter count.

\paragraph{Mathematical Problem Solving with Tool Use} We assess the capacity for mathematical reasoning supported by program generation on GSM8K and MATH datasets, leveraging few-shot program-of-thought prompting \citep{pot,pal}. The models are instructed to address each question by generating Python code, making use of libraries such as \textit{math} and \textit{sympy} to perform complex calculations as needed. The final solution is determined by executing the generated program and using its output as the answer. Table \ref{tab:base_math_pot_proof} illustrates that DeepSeekMath-Base 7B surpasses the previous best performance set by Llemma 34B.

\paragraph{Formal Mathematics}

Automating formal proofs offers significant advantages in verifying the correctness and dependability of mathematical arguments, while also boosting productivity—an area that has seen heightened research interest in recent years. In this study, we assess DeepSeekMath-Base 7B on the informal-to-formal proving benchmark described in \citep{dsp_proof}, which involves producing a formal proof by leveraging an informal statement, its corresponding formal version, and an informal proof. Our evaluation is conducted on the miniF2F dataset \citep{minif2f}, a standard testbed for formalizing challenging Olympiad-level mathematics, where the goal is to generate a formal Isabelle proof for each instance using few-shot prompting. Adopting the approach of \cite{dsp_proof}, we have the model construct proof outlines, and subsequently employ the widely used Sledgehammer automated prover \citep{sledgehammer} to complete the finer details. According to the outcomes in Table \ref{tab:base_math_pot_proof}, DeepSeekMath-Base 7B exhibits competitive capabilities in automated proof formalization.

\paragraph{Natural Language Understanding, Reasoning, and Code}

We assess the ability of the models in natural language understanding using MMLU \citep{mmlu}, reasoning with BBH \citep{bbh}, and code generation on HumanEval \citep{codex} and MBPP \citep{mbpp}. According to Table \ref{tab:understanding_reasoning_code}, DeepSeekMath-Base 7B demonstrates marked improvement over its predecessor, DeepSeek-Coder-Base-v1.5 \citep{deepseek-coder}, on both MMLU and BBH tasks, indicating that mathematical training bolsters both language understanding and reasoning skills. Furthermore, by incorporating code tokens during continual training, DeepSeekMath-Base 7B preserves the coding performance observed in DeepSeek-Coder-Base-v1.5 when evaluated on HumanEval and MBPP. In summary, DeepSeekMath-Base 7B achieves notably better results than the more generic Mistral 7B \citep{mistral} across all three benchmarks related to reasoning and coding.

\section{Supervised Fine-Tuning}

\subsection{SFT Data Curation}

We develop a mathematical instruction-tuning dataset that spans both English and Chinese tasks, encompassing a range of mathematical domains and degrees of difficulty. Each problem is matched with a solution articulated in chain-of-thought (CoT) \citep{cot}, program-of-thought (PoT) \citep{pot,pal}, or tool-integrated reasoning style \citep{tora}. The resulting dataset includes 776K training instances.

\begin{itemize}[topsep=0pt]
    \item \textbf{English mathematical datasets}: We provide tool-augmented annotations for GSM8K and MATH, in addition to incorporating a portion of MathInstruct \citep{MathInstruct} and the training set from Lila-OOD \citep{lila}, both of which feature solutions leveraging CoT or PoT reasoning strategies. The English dataset encompasses a broad spectrum of mathematical disciplines, including algebra, probability, number theory, calculus, and geometry.
    \item \textbf{Chinese mathematical datasets}: Our Chinese dataset comprises K-12 math questions across 76 sub-topics—such as linear equations—with detailed solutions presented in both CoT and tool-integrated formats.
\end{itemize}

\subsection{Training and Evaluating DeepSeekMath-Instruct 7B}

This section presents DeepSeekMath-Instruct 7B, a model obtained by applying mathematical instruction tuning to DeepSeekMath-Base. For training, we combine examples at random until the context window reaches up to 4K tokens. The model is trained over 500 steps, employing a batch size of 256 and maintaining a fixed learning rate of 5e-5.

To assess mathematical capabilities, we evaluate model performance in scenarios both lacking and incorporating external tool use, utilizing four quantitative reasoning benchmarks in English and Chinese. Our evaluation includes a comparison with state-of-the-art models available at the time:

\begin{itemize}[topsep=0pt]
    \item \textbf{Closed-source models} comprise the following: (1) the GPT series, with GPT-4 \citep{gpt4} and GPT-4 Code Interpreter~\footnote{\url{https://openai.com/blog/chatgpt-plugins\#\#code-interpreter}} representing the most advanced versions, (2) Gemini Ultra and Pro \citep{gemini}, (3) Inflection-2 \citep{inflection-2}, (4) Grok-1~\footnote{\url{https://x.ai/model-card}},
    as well as models made recently available by Chinese technology firms, including (5) Baichuan-3~\footnote{\url{https://www.baichuan-ai.com}},
    and (6) the most recent iteration of GLM-4~\footnote{\url{https://open.bigmodel.cn/dev/api\#glm-4}}

    from the GLM family \citep{glm}.

These models are primarily designed for broad applications and have typically passed through various alignment steps.

\item \textbf{Open-source models} feature: general-purpose systems such as (1) DeepSeek-LLM-Chat 67B \citep{deepseek-llm}, (2) Qwen 72B \citep{qwen}, (3) SeaLLM-v2 7B \citep{seallm}, and (4) ChatGLM3 6B \citep{chatglm3}. Models with a specific emphasis on mathematics include (5) InternLM2-Math 20B~\footnote{\url{https://github.com/InternLM/InternLM-Math}}, which extends InternLM2 with specialized mathematical pre-training and subsequent instruction tuning; (6) Math-Shepherd-Mistral 7B, which employs PPO training \citep{schulman2017proximal} on Mistral 7B \citep{mistral} using a process-supervised reward model; (7) the WizardMath series \citep{wizardmath}, which enhances the mathematical capabilities of Mistral 7B and Llama-2 70B \citep{llama2} via evolve-instruct (a type of instruction tuning utilizing AI-generated prompts) and PPO, mainly drawing training tasks from GSM8K and MATH; (8) MetaMath 70B \citep{metamath}, realized by fine-tuning Llama-2 70B with expanded datasets based on GSM8K and MATH; (9) ToRA 34B \cite{tora}, which is derived from CodeLlama 34B fine-tuned for integrating external tools in mathematical reasoning; and (10) MAmmoTH 70B \citep{MathInstruct}, comprising Llama-2 70B further trained with MathInstruct via instruction tuning.

As presented in Table \ref{tab:sft_rl_math}, when tool use is prohibited during evaluation, DeepSeekMath-Instruct 7B achieves robust step-by-step reasoning. Specifically, our model outperforms all open-source competitors and most closed-source models (such as Inflection-2 and Gemini Pro) on the competition-level MATH dataset by a margin of at least 9\% absolute. This advantage is maintained even over significantly larger models (for instance, Qwen 72B) and those optimized with math-centric reinforcement learning techniques (for example, WizardMath-v1.1 7B). While DeepSeekMath-Instruct is competitive with the Chinese proprietary systems GLM-4 and Baichuan-3 on MATH, it does not yet reach the performance level of GPT-4 or Gemini Ultra.

In a testing scenario that permits models to combine natural language reasoning with programmatic tool usage to tackle problems, DeepSeekMath-Instruct 7B achieves nearly 60\% accuracy on the MATH dataset, outperforming every open-source counterpart to date. On additional benchmarks, our model performs on par with DeepSeek-LLM-Chat 67B, the previously leading model despite being ten times larger in size.
\section{Reinforcement Learning}

\subsection{Group Relative Policy Optimization} Reinforcement learning (RL) has shown considerable success in enhancing the mathematical reasoning skills of large language models (LLMs) following the Supervised Fine-Tuning (SFT) phase \citep{wang2023math,wizardmath}. Here, we present our novel and efficient RL method, Group Relative Policy Optimization (GRPO).

\subsubsection{From PPO to GRPO} Proximal Policy Optimization (PPO) \citep{schulman2017proximal} constitutes an actor-critic reinforcement learning approach that has achieved broad adoption for the reinforcement learning fine-tuning (RLFT) phase in Large Language Models (LLMs) \citep{ouyang2022training}. This method updates LLMs by maximizing the following surrogate objective:

\begin{equation}
\footnotesize
    \mathcal{J}_{PPO}(\theta) = \mathbb{E}{[q \sim P(Q), o \sim \pi_{\theta_{old}}(O|q)]} \frac{1}{|o|} \sum_{t=1}^{|o|} \min \left[ \frac{\pi_\theta(o_{t} | q, o_{<t})}{\pi_{\theta_{old}}(o_{t} | q, o_{<t})} A_{t}, \text{clip} \left( \frac{\pi_\theta(o_{t} | q, o_{<t})}{\pi_{\theta_{old}}(o_{t} | q, o_{<t})}, 1 - \varepsilon, 1 + \varepsilon \right)  A_{t} \right] ,
\end{equation}

Here, $\pi_{\theta}$ denotes the current policy model, while $\pi_{\theta_{old}}$ signifies the previous iteration of the policy model. The variables $q$ and $o$ represent questions and the corresponding outputs, which are drawn from the dataset of questions and from samples generated by the prior policy $\pi_{\theta_{old}}$, respectively. The parameter $\varepsilon$ is a hyper-parameter used for clipping, incorporated within PPO to promote stability during training. The term $A_t$ corresponds to the advantage estimate, calculated using Generalized Advantage Estimation (GAE) \citep{gae}, informed by future rewards $\{r_{\ge t}\}$ and a value function $V_{\psi}$ trained for this purpose. Consequently, PPO entails the simultaneous training of a value function along with the policy. Furthermore, to address the potential issue of excessive reward optimization, it is conventional to introduce a KL penalty at each token step, using a reference model, as part of the per-token reward signal \citep{ouyang2022training}; specifically,

\begin{equation}
    r_{t} = r_\varphi(q, o_{\le t}) - \beta \log\frac{\pi_{\theta}(o_{t}|q, o_{<t})}{\pi_{ref}(o_{t}|q, o_{<t})},
\label{eq:PPO-reward}
\end{equation}

where $r_\varphi$  is the reward model, $\pi_{ref}$ is the reference model, which is usually the initial SFT model, and $\beta$  is the coefficient of the KL penalty.

The value function used in PPO is often as large and complex as the policy model itself, leading to notable computational and memory challenges. Furthermore, when training with reinforcement learning, the value function acts as a baseline to reduce the variance of the advantage estimate. In large language models (LLMs), however, rewards from the reward model are typically attributed solely to the final token, which complicates constructing a value function that reliably estimates returns at every token position. To tackle this issue, as illustrated in Figure \ref{fig:grpo}, we introduce Group Relative Policy Optimization (GRPO). This approach eliminates the need for the separate value function approximation used in PPO, and instead sets the baseline as the mean reward among several outputs generated for the same input question. Specifically, given each question $q$, GRPO generates a set of output samples $\{o_1, o_2, \cdots, o_G\}$ using the prior policy $\pi_{\theta_{old}}$, and then updates the policy model by optimizing the following objective:

\begin{equation}
\footnotesize
\begin{split}
    \mathcal{J}_{GRPO}(\theta) &= \mathbb{E}{[q \sim P(Q), \{o_i\}_{i=1}^G \sim \pi_{\theta_{old}}(O|q)]}  \\
    & \frac{1}{G}\sum_{i=1}^G\frac{1}{|o_i|} \sum_{t=1}^{|o_i|} \Bigg\{ \min \Bigg[ \frac{\pi_\theta(o_{i,t} | q, o_{i,<t})}{\pi_{\theta_{old}}(o_{i,t} | q, o_{i,<t})} \hat{A}_{i,t},\, \text{clip} \left( \frac{\pi_\theta(o_{i,t} | q, o_{i,<t})}{\pi_{\theta_{old}}(o_{i,t} | q, o_{i,<t})}, 1 - \varepsilon, 1 + \varepsilon \right)  \hat{A}_{i,t} \Bigg] - \beta\, \mathbb{D}_{KL}\left[\pi_{\theta} \,\Vert\, \pi_{ref}\right]\Bigg\} ,
\end{split}
\label{eq:GRPO-obj}
\end{equation}

where $\varepsilon$ and $\beta$ serve as hyper-parameters, and $\hat{A}_{i,t}$ denotes the advantage, which is derived by considering only the relative rewards among outputs within the same group—a procedure that will be explained further in the subsequent subsections. The methodology GRPO employs for advantage estimation, which is based on group-relative comparisons, is highly compatible with the inherently comparative setup of reward models, as these are usually trained using datasets consisting of paired output comparisons for identical questions. Furthermore, as opposed to incorporating a KL penalty directly into the reward, GRPO introduces regularization by explicitly appending the KL divergence between the learned policy and the reference policy to the loss function. This approach eliminates additional complexities in the computation of $\hat{A}_{i,t}$. Notably, distinct from the KL penalty utilized in (\ref{eq:PPO-reward}), the KL divergence here is approximated with the following unbiased estimator \citep{kl_approx}:

\begin{equation}
\small
    \mathbb{D}_{KL}\left[\pi_{\theta} || \pi_{ref}\right] = \frac{\pi_{ref}(o_{i,t}|q,o_{i,<t})}{\pi_{\theta}(o_{i,t}|q,o_{i,<t})}- \log\frac{\pi_{ref}(o_{i,t}|q,o_{i,<t})}{\pi_{\theta}(o_{i,t}|q,o_{i,<t})} - 1,
\end{equation}

which is guaranteed to be positive.

\begin{algorithm}[t]
  \small
  \caption{Iterative Group Relative Policy Optimization}
  \textbf{Input} Initial policy model $\pi_{\theta_{\text{init}}}$; reward models $r_\varphi$; set of task prompts $\mathcal{D}$; 
  hyperparameters $\varepsilon$, $\beta$, $\mu$
  \begin{algorithmic}[1]
    \State Set policy model $\pi_\theta \gets \pi_{\theta_{\text{init}}}$
    \For{iteration = 1, \dots, I}
       \State Assign reference model $\pi_{ref} \gets \pi_{\theta}$
      \For{step = 1, \dots, M}
      \State Draw a minibatch $\mathcal{D}_b$ from $\mathcal{D}$
      \State Update previous policy $\pi_{\theta_{old}} \gets \pi_{\theta}$ 
      \State For each prompt $q \in \mathcal{D}_b$, sample $G$ responses $\{o_i\}_{i=1}^G \sim \pi_{\theta_{old}} (\cdot \mid q)$
      \State Evaluate and assign reward scores $\{r_i\}_{i=1}^{G}$ for every sampled output $o_i$ using $r_{\varphi}$ 
      \State Estimate the group relative advantage $\hat{A}_{i,t}$ for the $t$-th token in each $o_i$
      \For{GRPO iteration = 1, \dots, $\mu$}
        \State Adjust the policy $\pi_{\theta}$ by optimizing the GRPO objective as defined in Equation \ref{eq:GC-GRPO}
      \EndFor
    \EndFor 
    \State Continuously update $r_\varphi$ via replay-based training.
    \EndFor 
  \end{algorithmic}
  \textbf{Output} $\pi_\theta$
  \label{alg:iter-grpo}
\end{algorithm}

\subsubsection{Outcome Supervision RL with GRPO} In this framework, given a question $q$, a collection of outputs $\{o_1, o_2, \cdots, o_G\}$ is generated by sampling from the previous policy $\pi_{\theta_{old}}$. Each output is then evaluated with a reward model, which assigns a set of $G$ reward scores $\mathbf{r}=\{r_1, r_2, \cdots, r_G\}$. These raw rewards undergo normalization: the mean reward of the group is subtracted from each score, and the result is divided by the group's standard deviation. Under outcome supervision, the normalized reward at the conclusion of output $o_i$ is used, assigning this value as the advantage $\hat{A}_{i, t}$ to all tokens comprising the output—specifically, $\hat{A}_{i, t} = \widetilde{r}_i = \frac{r_i- {\rm mean}(\mathbf{r})}{{\rm std}(\mathbf{r})}$. Policy optimization is performed by maximizing the objective outlined in equation (\ref{eq:GRPO-obj}).

\subsubsection{Process Supervision RL with GRPO} In outcome supervision, feedback is delivered only at the conclusion of the output, which may be inadequate or inefficient for guiding the learning process in intricate mathematical scenarios. To address this, in alignment with \cite{wang2023math}, we investigate process supervision, wherein rewards are provided after every individual reasoning step. Specifically, for a question $q$ and $G$ sampled responses $\{o_1, o_2, \cdots, o_G\}$, a dedicated process reward model evaluates and assigns scores to each step within the responses, resulting in rewards structured as: $\mathbf{R} = \{ \{r_1^{index(1)},\cdots,r_1^{index(K_1)}\}, \cdots,  \{r_G^{index(1)},\cdots,r_G^{index(K_G)}\} \}$. Here, $index(j)$ denotes the end token position of the $j$-th step, while $K_i$ designates the total number of steps in the $i$-th output. To promote consistency, rewards are standardized using their mean and standard deviation: $\widetilde{r}_i^{index(j)} = \frac{r_i^{index(j)} - {\rm mean(\mathbf{R})}}{{\rm std(\mathbf{R})}}$.

Next, within the process supervision framework, the advantage for each token is defined as the cumulative sum of the normalized rewards associated with all subsequent steps, expressed as $\hat{A}_{i, t} = \sum_{index(j) \ge t} \widetilde{r}_i^{index(j)}$. The policy is then updated by maximizing the objective given in equation (\ref{eq:GRPO-obj}).

\subsubsection{Iterative RL with GRPO} During the course of reinforcement learning, the previously trained reward model may become inadequate for effectively supervising the evolving policy model. To address this issue, we investigate an iterative RL approach employing GRPO. As illustrated in Algorithm \ref{alg:iter-grpo}, the iterative GRPO strategy involves constructing new datasets for the reward model by sampling trajectories from the latest policy model, followed by updating the reward model through a replay mechanism that retains 10\% of prior training samples. Subsequently, the updated policy model is assigned as the new reference model, and policy training proceeds in conjunction with the refreshed reward model.

\subsection{Training and Evaluating DeepSeekMath-RL}

Our reinforcement learning experiments are carried out using the DeepSeekMath-Instruct 7B model. For the RL phase, we utilize a dataset comprising chain-of-thought-formatted questions pertaining to GSM8K and MATH, sourced from the SFT data, totaling approximately 144K questions. To specifically assess the effect of RL on benchmarks devoid of corresponding data during RL, we intentionally exclude SFT questions from other domains. In constructing the reward model’s training dataset, we adhere to the methodology outlined in \citep{wang2023math}. The initial reward model is trained using DeepSeekMath-Base 7B with a learning rate of 2e-5. For policy optimization using GRPO, the policy model’s learning rate is set to 1e-6, and we adopt a KL divergence coefficient of 0.04. Each question prompts the sampling of $64$ responses, with a maximum sequence length limited to 1024 tokens and a batch size also fixed at 1024. The policy model undergoes only a single update after each exploration step. Evaluation of DeepSeekMath-RL 7B is conducted using the same suite of benchmarks as for DeepSeekMath-Instruct 7B. Specifically, for DeepSeekMath-RL 7B, the GSM8K and MATH benchmarks utilizing chain-of-thought reasoning are considered in-domain, while all remaining benchmarks are categorized as out-of-domain.

Table \ref{tab:sft_rl_math} presents the results of both open- and closed-source models employing chain-of-thought and tool-augmented reasoning strategies on English and Chinese evaluation sets. The analysis yields the following observations: 1) Using chain-of-thought reasoning, DeepSeekMath-RL 7B achieves accuracy rates of 88.2\% on GSM8K and 51.7\% on MATH. These outcomes not only exceed those recorded by every open-source model in the 7B to 70B range but also surpass the majority of closed-source systems. 2) Notably, DeepSeekMath-RL 7B’s training is limited to chain-of-thought style instruction tuning using GSM8K and MATH datasets, initializing from DeepSeekMath-Instruct 7B. In spite of this narrow training focus, it consistently outperforms DeepSeekMath-Instruct 7B across all metrics considered, highlighting the impact and advantages conferred by reinforcement learning.

\section{Discussion}
In this section, we will share our findings in pre-training and RL experiments. 

\subsection{Lessons Learnt in Pre-Training}

We begin by presenting insights gained from our pre-training phase. Unless stated otherwise, all experiments are conducted using the training configurations described in Section \ref{sec:quality-policy}. Notably, references to the DeepSeekMath Corpus within this section pertain to a dataset of 89B tokens, sourced from the second cycle of data collection.

\subsubsection{Code Training Benefits Mathematical Reasoning}

A widely cited, though unproven, hypothesis posits that exposure to code enhances reasoning abilities. Here, we provide evidence relevant to this claim in the context of mathematics: code training enhances models' proficiency in mathematical reasoning tasks, regardless of whether external tools are utilized.

To study how code training affects mathematical reasoning, we experimented with the following two-stage training and one-stage training settings:

\textbf{Two-Stage Training}

\begin{itemize}[topsep=0pt]
    \item \textbf{Code Training for 400B Tokens $\rightarrow$ Math Training for 150B Tokens}: We initially pretrain DeepSeek-LLM 1.3B using 400B code-related tokens, and subsequently continue its training with an additional 150B tokens drawn from mathematics datasets;
    \item \textbf{General Training for 400B Tokens $\rightarrow$ Math Training for 150B Tokens}: As a baseline, we carry out the same two-phase training, but in the first stage we substitute code tokens with general-domain tokens sourced from a broad corpus compiled by DeepSeek-AI, in order to assess whether pretraining on code confers benefits in mathematical reasoning compared to standard language data.
\end{itemize}

\textbf{One-Stage Training}

\begin{itemize}[topsep=0pt]
    \item \textbf{150B Token Mathematical Training}: We conduct training of DeepSeek-LLM 1.3B on a dataset consisting of 150B mathematical tokens;
    \item \textbf{Integrated Training with 400B Code Tokens and 150B Math Tokens}: Sequentially performing math training after code training is observed to reduce performance on coding tasks. Here, we explore whether jointly training on both code (400B tokens) and math (150B tokens) in a single stage can simultaneously enhance mathematical reasoning while mitigating the issue of catastrophic forgetting of code knowledge.
\end{itemize}

\paragraph{Results}
Table \ref{tab:code-math} and Table \ref{tab:code-math-general-eval} demonstrate the downstream performance under different training settings.

Code training enhances program-aided mathematical reasoning capabilities in both two-stage and one-stage training frameworks. As indicated in Table \ref{tab:code-math}, utilizing code training by itself during the initial stage of two-stage training already provides a substantial boost in solving GSM8K and MATH tasks using Python. Adding a subsequent math training phase yields additional performance gains. Notably, within the one-stage training paradigm, a combined mixture of code and math tokens not only counteracts the phenomenon of catastrophic forgetting observed with two-stage processes but also jointly improves both coding performance (Table \ref{tab:code-math-general-eval}) and program-assisted mathematical reasoning abilities (Table \ref{tab:code-math}).

Code training likewise enhances mathematical reasoning capabilities even in the absence of tool use. In a two-stage training paradigm, the preliminary phase focused on code already yields noticeable gains. Furthermore, this phase facilitates more effective math training in the second stage, culminating in optimal results. Conversely, employing a single-stage training regime that mixes both code tokens and math tokens tends to impair mathematical reasoning performance when tools are not available. One possible explanation is that DeepSeek-LLM 1.3B, given its relatively small model size, does not possess sufficient capacity to absorb and integrate code and mathematical content at the same time.

\subsubsection{ArXiv Papers Seem Ineffective in Improving Mathematical Reasoning}

ArXiv papers frequently constitute part of the math pre-training datasets \citep{minerva,gpt-f,llemma,mathpile}, yet their specific influence on enhancing mathematical reasoning remains largely unexplored. Somewhat unexpectedly, our experiments indicate that incorporating arXiv papers offers little benefit to mathematical reasoning performance. We evaluated this across models of varying capacities, such as DeepSeek-LLM 1.3B and DeepSeek-Coder-Base-v1.5 7B \citep{deepseek-coder}, using arXiv-based datasets processed through different pipeline configurations:

\begin{itemize}[topsep=0pt]
    \item \textbf{MathPile} \citep{mathpile}: comprising 8.9B tokens, this dataset is assembled through specialized cleaning and filtering approaches, with scientific arXiv papers making up more than 85\% of the collection;
    \item \textbf{ArXiv-RedPajama} \citep{redpajama}: consists of the full set of arXiv LaTeX files, in which preambles, comments, macros, and bibliographies have been excluded, resulting in a total of 28.0B tokens.
\end{itemize}

For our experimental setup, DeepSeek-LLM 1.3B was trained individually for 150B tokens, while DeepSeek-Coder-Base-v1.5 7B was trained for 40B tokens, using distinct arXiv corpora. Results indicate that incorporating arXiv papers does not significantly enhance models’ mathematical reasoning capabilities. When trained solely on arXiv data, both models demonstrate either negligible progress or even a decline in performance across an array of mathematical evaluation benchmarks utilized in this analysis. These benchmarks encompass quantitative reasoning datasets, such as GSM8K and MATH (Table \ref{tab:arxiv-cot}), multiple-choice tasks exemplified by MMLU-STEM (Table \ref{tab:arxiv-cot}), as well as formal mathematics tasks, including miniF2F (Table \ref{tab:arxiv_atp}).

However, this conclusion has its limitations and should be taken with a grain of salt. We have not yet studied:

\begin{itemize}[topsep=0pt]
    \item How arXiv tokens influence particular mathematical activities outside the scope of this study, for example, the informalization of theorems, which involves translating formal statements or proofs into their informal counterparts;
    \item The influence of arXiv tokens in conjunction with additional data sources;
    \item If scaling models up would amplify the advantages provided by incorporating arXiv papers.
\end{itemize}

Thus, further exploration is required, which we leave for future studies.

\subsection{Insights of Reinforcement Learning}
\subsubsection{Towards to a Unified Paradigm} 
In this part, we establish a comprehensive framework to systematically compare diverse training approaches, including SFT, RFT, DPO, PPO, and GRPO. We also present empirical studies to investigate the elements comprising this unified scheme. Typically, the gradient of any training objective with respect to parameter $\theta$ can be generalized as follows:

\begin{equation}
    \nabla_{\theta}\mathcal{J}_{\textcolor{red}{\mathcal{A}}}(\theta) = \mathbb{E}[\underbrace{(q,o) \sim \textcolor{red}{\mathcal{D}}}_{Data \ Source}]\left( \frac{1}{|o|} \sum_{t=1}^{|o|}  \underbrace{GC_{{\mathcal{A}}}(q, o, t, \textcolor{red}{\pi_{{rf}}})}_{Gradient \ Coefficient}  \nabla_{\theta}\log \pi_{\theta}(o_t | q, o_{<t})\right).
\label{eq:objective}
\end{equation}

The framework is composed of three essential elements: 1) the \textit{Data Source $\mathcal{D}$}, responsible for providing the training dataset; 2) the \textit{Reward Function $\pi_{{rf}}$}, which supplies the reward signals during training; and 3) the \textit{Algorithm $\mathcal{A}$}, which integrates the training data and reward signals to compute the gradient coefficient $GC$, thereby influencing the strength of applied penalties or rewards. Within this unified paradigm, we examine a number of representative techniques:

\begin{itemize}
    \item  \textbf{Supervised Fine-tuning (SFT)}: In SFT, the pretrained model undergoes further training on a dataset curated by human annotators specifically for fine-tuning purposes.
    \item \textbf{Rejection Sampling Fine-tuning (RFT)}: RFT builds upon the SFT model by additionally fine-tuning it on response samples, generated by the SFT model itself, that are screened based on their accuracy in answering SFT-derived questions.
    \item \textbf{Direct Preference Optimization (DPO)}: DPO enhances the SFT model through fine-tuning using a set of augmented response samples produced by the SFT model and optimized via a pairwise DPO loss function.
    \item \textbf{Online Rejection Sampling Fine-tuning (Online RFT)}: In contrast to standard RFT, Online RFT starts with the SFT model as the policy initialization and continues fine-tuning with dynamically sampled outputs from the current policy model.
    \item \textbf{PPO/GRPO}: PPO/GRPO procedures initialize the policy network with the SFT model and further adapt it through reinforcement learning with samples drawn on-the-fly from the policy.
\end{itemize}

We summarize the components of these methods in Table \ref{tab:data_policy}.
Please refer to Appendix \ref{app:analysis-rl} for a more detailed derivation process.

\paragraph{Observation about Data Source} We categorize data sources into two main types: online sampling and offline sampling. Online sampling refers to situations where training data is generated by the live exploration of the model being trained, whereas offline sampling involves using data collected from the initial SFT model’s outputs. Specifically, RFT and DPO utilize offline sampling, in contrast to Online RFT and GRPO, which employ online sampling.

As depicted in Figure \ref{fig:rl-analysis}, our results indicate that Online RFT achieves markedly better performance than RFT across both benchmarks. In the early training phase, Online RFT performs on par with RFT; however, it attains a clear performance edge as training progresses, highlighting the benefits of the online approach. This pattern is expected, given that the actor and SFT model are initially quite similar, resulting in only slight variances in the sampled data. As training advances, the actor-generated data increasingly diverge from the SFT model, making immediate and adaptive data sampling considerably more advantageous.

\paragraph{Observation about Gradient Coefficient} The reward signal is transformed into a gradient coefficient within the algorithm to facilitate model parameter updates. In our experimental setup, we decompose the reward function into `Rule' and `Model' components. The `Rule' approach evaluates a response's quality exclusively by checking the correctness of the answer, whereas the `Model' approach utilizes a trained reward model that assigns a score to each response. Notably, the reward model's training data are themselves generated based on rule-based assessments. As reflected in Equations \ref{eq:GC-RFT} and \ref{eq:GC-GRPO}, a primary distinction between GRPO and Online RFT emerges: GRPO dynamically modifies the gradient coefficient in accordance with the reward score outputted by the reward model. This design enables graded reinforcement and punishment tailored to the relative reward magnitude of each response. By comparison, Online RFT does not exhibit this flexibility; it applies equal positive reinforcement to all correctly answered responses and does not introduce negative reinforcement for incorrect ones.

As illustrated in Figure \ref{fig:rl-analysis}, GRPO outperforms online RFT, underscoring the advantage of modifying both positive and negative gradient coefficients. Moreover, GRPO+PS achieves better outcomes than GRPO+OS, suggesting that implementing step-specific, fine-grained gradient coefficients yields additional benefits. We also investigate iterative RL within our experimental framework by carrying out two iterations. As depicted in Figure \ref{fig:iter-rl}, iterative RL leads to marked performance gains, with the most pronounced improvement occurring after the initial iteration.

\subsubsection{Why RL Works?} In this study, we apply reinforcement learning using a selected subset of instruction tuning data, resulting in a substantial improvement over the baseline instruction-tuned model. To elucidate the mechanism underlying the effectiveness of reinforcement learning, we assess both Pass@K and Maj@K accuracies for the Instruct and RL-enhanced models across two evaluation benchmarks. Figure \ref{fig:maj-pass} demonstrates that RL consistently elevates Maj@K performance, yet does not impact Pass@K. These results suggest that reinforcement learning primarily strengthens the robustness of the model’s output distribution; that is, \textbf{the observed gains largely stem from increasing the likelihood of retrieving the correct answer from the TopK outputs, rather than advancing the model’s intrinsic reasoning abilities.} A comparable \textbf{misalignment problem} within the SFT model in reasoning contexts was documented by \citep{wang2023making}, who observed that implementing targeted preference alignment methods \citep{yuan2023rrhf,song2023preference,wang2023making} can notably enhance reasoning accuracy in SFT models.

\subsubsection{How to Achieve More Effective RL?}
\label{sec:effec_rl}
Our findings show that RL is quite effective for mathematical reasoning tasks. We further introduce a coherent framework that facilitates the interpretation of various prominent training approaches, treating each as a form of either standard or streamlined RL methodology. As outlined in Equation \ref{eq:objective}, the foundational elements include the Data Source, Algorithm, and Reward Function. Here, we suggest several avenues for advancing research on these three components.

\paragraph{Data Source} The data source serves as the foundational input for all training methodologies. Within Reinforcement Learning (RL), we define the data source as the set of unlabeled questions paired with responses generated by sampling from the policy model. For this study, our approach is limited to leveraging only those questions introduced during the instruction tuning phase, utilizing a straightforward nucleus sampling technique to produce outputs. We posit that this methodological constraint may explain why our RL procedure yields improvements exclusively in Maj@K performance. Moving forward, we aim to apply our RL framework to question prompts beyond the initial distribution and to integrate \textbf{advanced sampling (decoding) strategies}—such as those leveraging tree-search algorithms \citep{yao2023tree}. Additionally, recent \textbf{efficient inference techniques} \citep{xia-etal-2023-speculative,leviathan2023fast,kwon2023efficient,xia2024unlocking}, which critically influence the exploration effectiveness of policy models, are anticipated to be of substantial significance.

\paragraph{Algorithms} Algorithms utilize the data and reward signal to calculate the gradient coefficient, which in turn guides the update of model parameters. Drawing from Equation \ref{eq:objective}, contemporary approaches predominantly \textbf{TRUST} the reward function's signal to modulate the conditional probability of particular tokens, either increasing or decreasing it. Nevertheless, the reliability of the reward signal cannot always be guaranteed, particularly when dealing with highly complex tasks. For instance, even the extensively curated PRM800K datasets \citep{lightman2023let}, annotated by skilled professionals, still exhibit around 20\% of annotation errors\footnote{\url{https://github.com/openai/prm800k/issues/12\#issuecomment-1728491852}}. Consequently, we focus on investigating reinforcement learning algorithms capable of withstanding noisy reward signals. We argue that \textbf{WEAK-TO-STRONG} alignment strategies \citep{burns2023weak} have the potential to fundamentally transform learning algorithms.

\paragraph{Reward Function} The reward function provides the essential training signal in reinforcement learning. Typically, this function is implemented as a neural reward model. We identify three key areas for further development of reward models: 1) \textbf{Improving the generalization capacity of reward models.} It is crucial for the reward model to perform well on questions that lie outside of its training distribution, as well as on more complex decoding results. Failing to do so may result in reinforcement learning only reinforcing the existing output distribution of LLMs, without actually enhancing their capabilities; 2) \textbf{Incorporating uncertainty estimation in reward models.} Capturing and quantifying the reward model’s uncertainty could facilitate the interaction between weak reward models and algorithms designed for weak-to-strong learning; 3) \textbf{Developing high-quality process reward models efficiently} so that they yield nuanced feedback during the reasoning process, enabling more fine-grained learning signals \citep{lightman2023let,wang2023math}.

\section{Conclusion, Limitation, and Future Work} 

In this paper, we introduce DeepSeekMath, which establishes new state-of-the-art performance among open-source models on the competition-level MATH benchmark and nearly matches the results of proprietary systems. DeepSeekMath~is based on DeepSeek-Coder-v1.5 7B and is further trained with 500B additional tokens, including a substantial 120B math-related tokens predominantly extracted from Common Crawl. Through a comprehensive ablation analysis, we find that web-sourced pages can serve as valuable sources of mathematical data, whereas arXiv content does not yield the anticipated gains. Additionally, we propose Group Relative Policy Optimization (GRPO), a modified version of Proximal Policy Optimization (PPO) that enhances mathematical reasoning performance with reduced memory usage. Experimental findings confirm that GRPO remains beneficial even when DeepSeekMath-Instruct 7B already attains strong benchmark results. Lastly, we put forward a cohesive framework to interpret a range of related techniques and outline several promising avenues for advancing reinforcement learning efficacy.

While DeepSeekMath~demonstrates strong performance on quantitative reasoning tasks, its abilities in geometry and theorem-proving tasks remain inferior to those of closed models. As observed in our preliminary testing, the model struggles with triangle and ellipse-related questions, suggesting the possibility of a pre-training and fine-tuning data bias. Moreover, DeepSeekMath~is constrained by its model size, resulting in inferior few-shot learning compared to GPT-4. Whereas GPT-4 benefits from improved accuracy with the inclusion of few-shot examples, DeepSeekMath~performs comparably in both zero-shot and few-shot settings. Looking ahead, we aim to refine our engineered data selection methods to assemble a higher-quality pre-training dataset. We also plan to investigate additional strategies (see Section \ref{sec:effec_rl}) for enhancing reinforcement learning techniques in LLMs.

\end{document}